\begin{table}[H]
\centering
\caption{两套赛程的可观测结构与规模强度比较}
\label{tab:q4-structural-comparison}
\small
\resizebox{\textwidth}{!}{%
\begin{tabular}{lllrrl}
\toprule
指标 & 定义与分母 & 单位 & 2022实际 & 2026优化 & 偏好方向 \\
\midrule
最小休息 & 球队相邻比赛 UTC 间隔的最小值 & 小时 & 87.0000 & 63.0000 & 越高越优 \\
平均休息 & 球队相邻比赛 UTC 间隔的平均值 & 小时 & 97.8125 & 139.8125 & 越高越优 \\
跨时区次数 & 相邻比赛场馆 UTC 偏移发生变化的次数 & 次 & 0 & 78 & 越低越优 \\
跨时区率 & 跨时区次数除以球队转场次数 & 比例 & 0.0000 & 81.2500 & 越低越优 \\
场馆旅行 & 相邻比赛实际场馆大圆距离的平均值 & 公里/次 & 18.9868 & 2019.1075 & 越低越优 \\
换馆率 & 相邻轮更换场馆的球队转场比例 & 比例 & 78.1250 & 94.7917 & 越低越优 \\
场馆日利用率 & 发生比赛的场馆日占可用场馆日比例 & 比例 & 45.1923 & 25.0000 & 越高越优 \\
场馆日峰值 & 单一场馆单日承办场次最大值 & 场/场馆日 & 2 & 1 & 越低越优 \\
容量匹配 & 预计现场观众与场馆容量之比的场均值 & 比例 & 81.7689 & 61.7245 & 越高越优 \\
黄金覆盖 & 按赛事进度日和 UTC 钟点映射后落入前四分位时段的比例 & 比例 & 25.0000 & 62.5000 & 越高越优 \\
场均预计观众 & 预计现场观众总量除以比赛场次 & 人/场 & 41451.67 & 41704.92 & 越高越优 \\
\bottomrule
\end{tabular}
}
\vspace{2pt}
\noindent\begin{minipage}{0.96\textwidth}\footnotesize 注：休息、时区、场馆路径与场馆日由逐场赛程直接复算；黄金覆盖来自唯一时段映射，容量匹配和预计观众的需求端为题内赛前代理。实际赛程有64次球队转场，优化赛程有96次。\end{minipage}
\end{table}
